\documentclass[
    language=english,
    doctype=lecture-notes,
    institution=none,
]{../../omnilatex}

\RequirePackage{config/document-settings}

\begin{document}
\maketitle

\section{Eigenvalues and Eigenvectors}

\subsection{Definitions and Motivation}

We begin with the central definition of this lecture.

\begin{definition}
Let $A$ be an $n \times n$ matrix over $\mathbb{R}$.  A nonzero vector
$\mathbf{v} \in \mathbb{R}^n$ is an \textbf{eigenvector} of $A$ if there
exists a scalar $\lambda \in \mathbb{R}$ such that
\[
    A\mathbf{v} = \lambda \mathbf{v}.
\]
The scalar $\lambda$ is called the corresponding \textbf{eigenvalue}.
\end{definition}

\begin{remark}
Geometrically, an eigenvector is a direction that is preserved (up to scaling)
under the linear transformation defined by $A$.  The eigenvalue tells us the
factor by which the eigenvector is stretched or compressed.
\end{remark}

\subsection{The Characteristic Polynomial}

\begin{definition}
The \textbf{characteristic polynomial} of an $n \times n$ matrix $A$ is
\[
    p_A(\lambda) = \det(A - \lambda I).
\]
The eigenvalues of $A$ are precisely the roots of $p_A(\lambda) = 0$.
\end{definition}

\begin{theorem}
An $n \times n$ matrix has at most $n$ eigenvalues (counted with algebraic
multiplicity).
\end{theorem}

\begin{proof}
Since $\deg(p_A) = n$, the Fundamental Theorem of Algebra guarantees at most
$n$ roots.
\end{proof}

\subsection{Computing Eigenvalues}

\begin{example}
Find the eigenvalues of
\[
    A = \begin{pmatrix} 4 & 1 \\ 2 & 3 \end{pmatrix}.
\]

\textbf{Solution.}  We compute:
\[
    \det(A - \lambda I) = \det\begin{pmatrix} 4-\lambda & 1 \\ 2 & 3-\lambda \end{pmatrix}
    = (4-\lambda)(3-\lambda) - 2 = \lambda^2 - 7\lambda + 10
    = (\lambda - 5)(\lambda - 2).
\]
The eigenvalues are $\lambda_1 = 5$ and $\lambda_2 = 2$.
\end{example}

\begin{example}
Find the eigenvectors of the matrix $A$ from the previous example.

\textbf{Solution.}  For $\lambda_1 = 5$:
\[
    (A - 5I)\mathbf{v} = \begin{pmatrix} -1 & 1 \\ 2 & -2 \end{pmatrix}\mathbf{v} = \mathbf{0}
    \implies \mathbf{v}_1 = \begin{pmatrix} 1 \\ 1 \end{pmatrix}.
\]
For $\lambda_2 = 2$:
\[
    (A - 2I)\mathbf{v} = \begin{pmatrix} 2 & 1 \\ 2 & 1 \end{pmatrix}\mathbf{v} = \mathbf{0}
    \implies \mathbf{v}_2 = \begin{pmatrix} 1 \\ -2 \end{pmatrix}.
\]
\end{example}

\section{Diagonalization}

\begin{definition}
An $n \times n$ matrix $A$ is \textbf{diagonalizable} if there exists an
invertible matrix $P$ and a diagonal matrix $D$ such that $A = PDP^{-1}$.
\end{definition}

\begin{theorem}
An $n \times n$ matrix is diagonalizable if and only if it has $n$ linearly
independent eigenvectors.
\end{theorem}

\begin{proof}
Suppose $A$ has $n$ linearly independent eigenvectors $\mathbf{v}_1,
\ldots, \mathbf{v}_n$ with eigenvalues $\lambda_1, \ldots, \lambda_n$.
Let $P = [\mathbf{v}_1 \mid \cdots \mid \mathbf{v}_n]$ and $D =
\operatorname{diag}(\lambda_1, \ldots, \lambda_n)$.  Then $AP = PD$,
so $A = PDP^{-1}$ since $P$ is invertible.  The converse follows by
reversing the construction.
\end{proof}

\begin{example}
Diagonalize $A = \begin{pmatrix} 4 & 1 \\ 2 & 3 \end{pmatrix}$.

\textbf{Solution.}  Using the eigenvectors from the previous example:
\[
    P = \begin{pmatrix} 1 & 1 \\ 1 & -2 \end{pmatrix}, \quad
    D = \begin{pmatrix} 5 & 0 \\ 0 & 2 \end{pmatrix}, \quad
    P^{-1} = \frac{1}{3}\begin{pmatrix} 2 & 1 \\ 1 & -1 \end{pmatrix}.
\]
One can verify $A = PDP^{-1}$.
\end{example}

\section{Key Concepts}

\begin{description}
    \item[Eigenvalue] A scalar $\lambda$ such that $A\mathbf{v} = \lambda\mathbf{v}$
          for some nonzero vector $\mathbf{v}$.
    \item[Eigenvector] A nonzero vector $\mathbf{v}$ satisfying $A\mathbf{v} =
          \lambda\mathbf{v}$.
    \item[Characteristic polynomial] $p_A(\lambda) = \det(A - \lambda I)$;
          its roots are the eigenvalues.
    \item[Diagonalization] Expressing $A$ as $PDP^{-1}$ where $D$ is diagonal.
    \item[Algebraic multiplicity] The multiplicity of $\lambda$ as a root of
          $p_A$.
    \item[Geometric multiplicity] The dimension of the eigenspace corresponding
          to $\lambda$.
\end{description}

\section{Practice Problems}

\begin{problem}
Find the eigenvalues and eigenvectors of
\[
    B = \begin{pmatrix} 2 & 0 & 0 \\ 0 & 3 & 4 \\ 0 & 4 & 9 \end{pmatrix}.
\]
\end{problem}

\begin{problem}
Determine whether the matrix
\[
    C = \begin{pmatrix} 1 & 1 \\ 0 & 1 \end{pmatrix}
\]
is diagonalizable.  Justify your answer.
\end{problem}

\begin{problem}
Prove that a matrix $A$ is singular if and only if $0$ is an eigenvalue of $A$.
\end{problem}

\section{Summary}

Today we introduced eigenvalues and eigenvectors, which capture the fundamental
geometric behavior of a linear transformation.  The characteristic polynomial
provides a computational tool for finding eigenvalues, and diagonalization
allows us to decompose a matrix into a form that makes powers and matrix
exponentials easy to compute.  These concepts will be essential for our upcoming
discussion of quadratic forms and singular value decomposition.

\noindent\textbf{Next lecture:} Singular Value Decomposition and its applications.

\end{document}
